% ==========================================
% SECCIÓN 4: PROGRAMA
% ==========================================
% TIP: Esta sección describe el programa desarrollado para la práctica, incluyendo su estructura, funcionalidades, y cómo se implementaron los conceptos de tipos de datos, procesos y archivos en C.
%
% TIP: Las rutas de imágenes son relativas a main.tex.
% ==========================================
% RECORDATORIO DE REQUISITOS:
% ==========================================
% - Se debe mostrar la compilación en el reporte y esta no debe tener errores o advertencias
% - A continuación se debe mostrar la ejecución de los programas
% - Si no funciona el sistema o programa no se califica la práctica
% - No se pueden utilizar sockets para comunicar procesos
% - No se pueden utilizar gestores de bases de datos para almacenar la información de usuarios o productos
% - Se deben validar todos los campos donde se introduzcan datos por parte del usuario o cliente
% - Se deben subir todos los archivos .c y .h generados
% - No deben archivar o comprimir los archivos .c y .h
% - Solo se puede utilizar lenguaje C para la implementación de los programas

% ===============================================================

\chapter{Desarrollo}
\label{cap:programa}

% ------------------------------------------
% Responsable: Bustillos Cruz Jonatan
% ------------------------------------------

% TODO (Jonatan): Describir el programa desarrollado para la práctica, incluyendo su estructura, funcionalidades, y cómo se implementaron los conceptos de tipos de datos, procesos y archivos en C. Incluir ejemplos de código relevantes y capturas de pantalla de la ejecución del programa. Asegurarse de mostrar la compilación sin errores ni advertencias, y validar todos los campos donde se introduzcan datos por parte del usuario o cliente.


